A set of contributions in learning-based navigation relate to this work, alongside efforts to enable avoidance of thin obstacles. Many works have employed deep learning techniques to tackle the issue of autonomous robot navigation, with a particular focus on using RGB/depth cameras. These sensors have garnered increased attention due to their low cost, low power consumption, and lightweight design.
The authors in~\cite{Loquercio2021Science,Tolani2021visual} use imitation learning to generate collision-free smooth trajectories from realistic simulated data.
Though much effort has been devoted to obtaining more realistic simulation images~\cite{Zakharov2022photo_realistic,Bousmalis2018DomainAdaptation}, there persist effects that are difficult to model such as missing depth pixels when observing thin obstacles or low-texture, shiny surfaces.
In~\cite{Fereshteh2017CAD2RL}, reactive navigation policies are learned through reinforcement learning and domain randomization using RGB image inputs.
On the other hand,~\cite{Gandhi2017learning2fly,Kahn2021BADGR,Kahn2021LAND} predict probabilities of events such as collisions, going over bumpy terrain, or human disengagement from real-world RGB data.
Our approach also leverages real depth data in the training pipeline, however, part of the depth frames and the data collection where the robot dynamics are rolled out and collision events happen is performed in simulation.
The authors in~\cite{Hoeller2021representation} propose another method to mix real and simulated depth images for learning an environment representation for legged robot navigation. In this work, we focus on using both the real and simulated depth images augmented with the labeled semantic masks to ensure the reconstruction of hard-to-perceive obstacles such as thin objects in the environment representation.

Of further relevance to this work are also contributions proposing solutions to allow robots to detect and avoid thin obstacles in the environments. The authors in~\cite{Zhou2017FastAT} present a thin-structure obstacle detection and 3D reconstruction method based on an edge detector and edge-based visual odometry techniques. On the other hand, a multi-view algorithm for wire reconstruction using a parametric catenary model is illustrated in~\cite{Madaan2019MultiviewRO}. The work in~\cite{Landry2016MIP} proposes a planning approach using mixed integer programming to derive collision-free trajectories for a quadrotor to fly in environments populated with thin strings, however, the locations of the obstacles are given beforehand as convex hulls. The authors in~\cite{Dubey2018droan} use a monocular wire detection method based on synthetic data and a dilated convolutional neural network (CNN)~\cite{Madaan2017WireDU} to perceive thin-wire obstacles. The obstacles are then represented by a disparity-space representation for collision checking. A separate planning module evaluates a pre-computed trajectory library to choose and execute the best trajectory.
